\Sec{Conclusão}

\Subsec{Considerações sobre o trabalho desenvolvido}

Ao longo do trabalho, analisámos o funcionamento interno de uma ALU (\textit{Arithmetic Logic Unit}), a sua estrutura e o modo como esta realiza operações aritméticas e lógicas — funções estas que estão na base de qualquer aparelho/dispositivo digital que lide com cálculos e decisões lógicas.

\vspace{5mm}

A ALU, afinal, reflete a interação entre a teoria da matemática e a lógica digital, que em conjunto, permitem a construção de sistemas capazes de realizar tarefas que variam desde operações simples até às mais complexas.

\vspace{5mm}

Em suma, ao estudar e entender, logica e objetivamente, uma ALU, percebemos como funciona cada módulo desta, e adquirindo como bases, alguns módulos de trabalhos anteriores, pudemos aplicar e organizar a forma e maneira como se comportavam as entradas e saídas de cada parte da nossa ALU, a modo de atingir o objetivo proposto pelo enunciado.



\Subsec{Discussão de Resultados}

No que diz respeito aos resultados obtidos neste segundo trabalho, podemos afirmar que houve uma conclusão de resultados efetiva e que está de acordo com o esperado. Contudo, durante a realização deste projeto, notamos que, devido à complexidade da criação de cada módulo desta ALU, é provável que surjam algumas irregularidades no que diz respeito ao código VHDL. Mesmo assim, conseguimos ultrapassar cada falha que foi surgindo, alcançando assim, um módulo da ALU conforme previsto inicialmente.

\vspace{5mm}

Essa unidade é bastante importante, uma vez que possibilita a resolução de problemas e operações envolvendo lógica e aritmética.

\vspace{5mm}

De modo geral, foi essencial estabelecermos um conhecimento inicial sobre uma ALU para garantir uma melhor orientação e consistência ao longo de todo o trabalho. Além disso, a verificação de cada caso particular de uma operação, deve ser muito bem organizada, especialmente com o auxílio de uma \textit{testbench} previamente configurada e bem estruturada, conforme o código.



